
This appendix develops the categorical and algebraic core of the
FRFP (Foundational Reasoning and Feedback Protocol) architecture.
It provides a self-contained account of:

\begin{itemize}
  \item an ambient category with explicit and tacit modes;
  \item a free explicit algebra of task decompositions (the Task Decomposition Grammar, TDG);
  \item a navigation groupoid acting on explicit shapes;
  \item a fibration of explicit pipelines over shapes;
  \item a Grothendieck-style semidirect product \(N \ltimes E\);
  \item a semantic functor into a preorder-enriched tacit category; and
  \item correctness and canonicality theorems for Epistemic Structure.
\end{itemize}

We assume familiarity with basic category theory, rewriting systems,
and context-free grammars.

Throughout, we write \(\Cfull\) for the ambient category,
\(\Espace\) for explicit space, and \(\Tspace\) for tacit space.

%======================================================================
\subsection{Foundations}
%======================================================================

\subsubsection{Ambient Category \texorpdfstring{\(\Cfull\)}{Cfull}}

\begin{definition}[Ambient category]
The ambient category \(C_{\mathrm{full}}\) is an arbitrary category that may contain objects beyond those used by the FRFP protocol. We designate three \emph{distinguished} objects
\(\varnothing, E, T \in \mathrm{Ob}(C_{\mathrm{full}})\):
\begin{itemize}
  \item \(\varnothing\): the empty explicit state;
  \item \(E\): the explicit state object (machine-manipulable representations);
  \item \(T\): the tacit state object (human correctness-bearing states).
\end{itemize}
Morphisms in \(\Cfull\) are partitioned into three epistemic classes:
\begin{itemize}
  \item explicit morphisms with codomain \(E\), forming the explicit subcategory
        \(\Espace \subseteq \Cfull\);
  \item tacit morphisms with codomain \(T\), forming the tacit subcategory
        \(\Tspace \subseteq \Cfull\);
  \item boundary morphisms \(E \to T\), of which a distinguished one is the
        \emph{Representational Backflow} operator
        \(\RB : E \to T\).
\end{itemize}
Composition in \(\Cfull\) satisfies the usual category axioms:
\[
(h \circ g) \circ f \;=\; h \circ (g \circ f), \qquad
f \circ \id_X \;=\; f \;=\; \id_Y \circ f.
\]
\end{definition}

\begin{definition}[Explicit--Tacit Separation (ETS)]
The \emph{Explicit--Tacit Separation} axiom requires:
\begin{itemize}
  \item there are no morphisms \(T \to E\) in \(\Cfull\);
  \item tacit evaluation cannot be embedded inside explicit pipelines;
  \item the boundary morphism \(\RB : E \to T\) is not invertible.
\end{itemize}
\end{definition}

\subsubsection{Modal Interpretation of Explicit and Tacit Knowledge}

\begin{definition}[Explicit modality]
A cognitive item \(k\) is represented in \emph{explicit} modality if there exist
maps
\[
\enc : E \to \Sigma^{\ast}, \qquad
\dec : \Sigma^{\ast} \to E
\]
such that \(\dec(\enc(k))\) is equivalent to \(k\) in its epistemic role
(e.g., as a proof, equation, or program). In particular, \(k\) can be
faithfully encoded as a finite symbolic artifact and recovered without
loss of its correctness-relevant structure.
\end{definition}

\begin{definition}[Tacit modality]
A cognitive item is represented in \emph{tacit} modality when no such
encoding preserves its correctness-bearing content. Judgments, intuitions,
acceptance decisions, and other human correctness-bearing acts live in
tacit space.
\end{definition}

Human expression from tacit cognition to explicit content (e.g., writing a
new question, revising a claim) is modeled externally as the appearance of
new explicit morphisms in \(\Espace\); it is not modeled as a morphism
\(T \to E\).

\subsubsection{Explicit Category \texorpdfstring{\(\Espace\)}{E}}

\begin{definition}[Explicit category]
The explicit category \(\Espace\) is the full subcategory of \(\Cfull\)
on objects with codomain \(E\). Its core generating morphisms are:
\[
\RI : \varnothing \to E
\quad\text{(Representation Initiation)},\qquad
\EC : E \to E
\quad\text{(Explicit Computation)},\qquad
\ED : E \to E
\quad\text{(Explicit Diagnostics)}.
\]
\end{definition}

\noindent
\textbf{Remark.}
\(\EC\) represents algorithmic transformations of explicit content
(e.g., inference, rewriting, search). \(\ED\) represents explicit
checks and diagnostics (e.g., consistency checks, sanity tests),
preparing but not performing tacit evaluation.

\subsubsection{Tacit Category \texorpdfstring{\(\Tspace\)}{T}}

We now formalize tacit space as a preorder-enriched one-object category.

\begin{definition}[Tacit category]
The tacit category \(\Tspace\) has:
\begin{itemize}
  \item a single object, denoted again by \(T\);
  \item hom-set \(\Hom_{\Tspace}(T,T)\) equipped with a composition
        operation \(\circ\), making it a monoid;
  \item a preorder \(\sqsubseteq\) on \(\Hom_{\Tspace}(T,T)\) such that
        composition is monotone:
        \[
        f_1 \sqsubseteq f_2,\ g_1 \sqsubseteq g_2
        \quad\Rightarrow\quad
        g_1 \circ f_1 \sqsubseteq g_2 \circ f_2.
        \]
\end{itemize}
\end{definition}

\begin{definition}[Tacit primitives]
The tacit primitives are the morphisms
\[
\TE : T \to T \quad\text{(Tacit Evaluation)},\qquad
\HFD : T \to T \quad\text{(Human Final Decision)},
\]
subject to:
\begin{itemize}
  \item \emph{idempotence}:
        \(\TE \circ \TE = \TE\) and \(\HFD \circ \HFD = \HFD\);
  \item \emph{commutation}:
        \(\TE \circ \HFD = \HFD \circ \TE\);
  \item \emph{monotonicity}: for every \(f : T \to T\),
        \[
        f \sqsubseteq \TE \circ f,\qquad
        f \sqsubseteq \HFD \circ f.
        \]
\end{itemize}
\end{definition}

\begin{definition}[Correctness quotient]
Let \(J\) be a finite set of correctness judgments
(e.g., \(\{\mathsf{accept}, \mathsf{reject}, \mathsf{revise}\}\)).
A \emph{correctness quotient} is a surjective map
\[
\gamma : \Obj(\Tspace) \to J
\]
assigning to each tacit state a correctness class, subject to
\[
\gamma \circ \TE = \gamma = \gamma \circ \HFD.
\]
Thus tacit evaluation and final decision may refine or stabilize a tacit
state but do not change its correctness class.
\end{definition}

\paragraph{Interpretation and non-triviality of \texorpdfstring{$T$}{T}.}
The tacit category T is not introduced merely as a formal firewall that “reserves correctness to tacit space”.
Rather, it is intended to capture a substantive empirical claim about long-horizon scientific and regulatory work:
in these regimes, there are correctness-bearing aspects of judgment that cannot be exhaustively encoded in a fixed,
finite explicit representation without loss.

Definitions A.3–A.4 distinguish explicit and tacit modalities by the existence of lossless encodings into symbolic artefacts.
In many familiar settings—for example, textbook exercises, small programs, or contest-style benchmarks—it is reasonable
to treat \emph{inner} task correctness as effectively explicit: the task is specified by a finite prompt, there is a fixed
input–output relation or reference implementation, and success can be checked by a test harness or proof checker.
FRFP does not deny the existence or value of such regimes. Instead, it treats them as \emph{tacit-insensitive inner layers}:
conditional on a surrounding tacit backdrop (choice of formalism, modelling assumptions, risk tolerances),
correctness for that inner task behaves as if it were a pure explicit oracle on E.

The global explicit/tacit split remains.
Even in proof-assistant or test-oracle settings, tacit judgment is involved in selecting the calculus, specifying what the
theorem or contract should state, and deciding whether conformance to that specification is adequate for the real-world
purpose.
What FRFP targets are workflows where this outer tacit layer is non-trivial and dynamically involved in long-horizon
use: tasks evolve, standards and constraints shift, and background scientific understanding and norms play an essential
role in deciding whether an explicit artefact is “good enough”.
In those settings it is non-trivial, and in our view appropriate, to treat the final correctness status of an explicit
pipeline as living in T rather than in E. (See Appendix~D.4 for further discussion.)


A reasonable skeptic might ask: under what conditions could one reject the Epistemic Spaces axioms
and the later Epistemic Algebra structure on $T$? There are at least three salient counter-scenarios:
\begin{enumerate}
  \item \textbf{Tacit-insensitive inner correctness.}
For some domains and layers, once a formalism, specification, and world-modelling stance have been tacitly fixed,
every correctness-bearing judgment at the inner level can be captured by a finite explicit specification and a sound
and complete proof- or test-based oracle.
At that inner layer one can model correctness directly in E and treat T as abstracted away; FRFP’s Phase–3
inevitability and non-automation results then become vacuous or uninteresting for that particular subproblem.
Our focus is on regimes where the outer tacit layer—where specifications, standards, and modelling choices are
themselves evolving—remains dynamically involved in correctness.

  \item \textbf{No degradation or escalation.} Epistemic Algebra impossibility results rely on a
  degradation operator $\delta$ and recurring nontrivial requirements relative to a floor $r$.
  If one believes that, in their target setting, tacit state does not degrade in any relevant
  sense, or that requirements do not escalate or recur in the way we assume (for example,
  because tasks are short-lived and fully specified at the outset), then the inevitability
  theorems may not apply. Such systems fall outside FRFP's intended long-horizon regime.
  \item \textbf{Perfect tacit emulation.} If one posits an AI system that can internally
  reconstruct or simulate human tacit evaluation from explicit traces without loss---violating
  the NTER and CSC constraints---then our explicit--tacit separation (ETS, HEG, HEC) is simply
  rejected. FRFP explicitly rules out this possibility by axiom rather than attempting to model
  it.
\end{enumerate}

In summary, treating correctness as localized in $T$ is a substantive modelling choice:
it reflects the view that, for many real scientific and governance workflows, there is residual
human judgment that cannot be fully captured by static benchmarks or formal specifications, and
that this residual plays an essential role in long-horizon correctness. The impossibility and
inevitability results of Epistemic Structure and Dynamics should therefore be read as conditional on this non-trivial
tacit structure: they apply precisely in domains where correctness genuinely outstrips what can
be made explicit and mechanically checked.


\subsubsection{Boundary Morphism and Semantic RB Functor}

\begin{definition}[Boundary morphism \(\RB\)]
The \emph{Representational Backflow} operator is a morphism
\[
\RB : E \to T
\]
in \(\Cfull\). It is the unique explicit-to-tacit interface specified by
the Representational Backflow axiom.
\end{definition}

Let \(E_{\mathrm{sem}}\) denote the free explicit category generated by the
Task Decomposition Grammar (defined in Section~\ref{sec:appendix-tdg}).

\begin{definition}[Semantic RB functor]
The \emph{semantic RB} is a functor
\[
\RBsharp : E_{\mathrm{sem}} \to \Tspace
\]
assigning to each explicit pipeline a tacit state representing the content
received by the human for evaluation.
\end{definition}

\subsubsection{Interaction Protocols and Epistemic Spaces Axioms}

\begin{definition}[Interaction Protocols constraints]
The operational environment satisfies:
\begin{itemize}
  \item \textbf{IL (Interaction Limitation):} interaction is bounded and turn-based;
  \item \textbf{AR (Alignment Restriction):} the system must not autonomously act
        outside explicit instructions;
  \item \textbf{MD (Mode Discipline):} separation between explicit and tacit modes
        is enforced at the interface level;
  \item \textbf{NTER (No Tacit Evaluation Reconstruction):} the system must not
        reconstruct tacit evaluation internally from explicit traces;
  \item \textbf{CSC (Correctness Separation Constraint):} correctness decisions
        must not be delegated to explicit computation.
\end{itemize}
\end{definition}

\begin{definition}[Epistemic Spaces axioms]
An admissible epistemic framework satisfies:
\begin{itemize}
  \item \textbf{ETS (Explicit--Tacit Separation):} no morphisms \(T \to E\);
  \item \textbf{HEG (Human-Exclusive Grounding):} correctness grounding occurs
        only in tacit space via \(\TE\), executed by humans;
  \item \textbf{HEC (Human-Exclusive Closure):} final acceptance/rejection occurs
        only via \(\HFD\), executed by humans;
  \item \textbf{RB (Representational Backflow):} there is a unique boundary
        morphism \(\RB : E \to T\) used by all pipelines;
  \item \textbf{CP (Compositional Pipelines):} epistemic processes are built as
        compositions of primitives into pipelines;
  \item \textbf{AE (AI-Explicit Restriction):} AI systems act only via
        explicit-space morphisms in \(\Espace\).
\end{itemize}
\end{definition}

%======================================================================
\subsection{Explicit Algebra and the Task Decomposition Grammar}
\label{sec:appendix-tdg}
%======================================================================

\subsubsection{TDG Signature and Free Category}

\begin{definition}[TDG signature]
The Task Decomposition Grammar (TDG) signature \(\Sigma_{\mathrm{TDG}}\)
consists of:
\begin{itemize}
  \item objects: \(\varnothing\) and \(E\);
  \item generating morphisms:
  \[
    \RI : \varnothing \to E,\qquad
    \EC : E \to E,\qquad
    \ED : E \to E;
  \]
  \item structural operators \(\BRANCH\), \(\PARALLEL\), \(\LOOP\) acting
        as higher-arity constructors on morphisms.
\end{itemize}
\end{definition}

\begin{definition}[Free explicit category]
The explicit semantic category is the free category on the TDG signature:
\[
E_{\mathrm{sem}} := \FreeCat(\Sigma_{\mathrm{TDG}}).
\]
There is a canonical interpretation functor
\(\Ufun : E_{\mathrm{sem}} \to \Espace\)
sending each TDG term to its denotation as an explicit morphism.
\end{definition}

\begin{theorem}[Free explicit algebra]
\label{thm:free-explicit-algebra}
For any category \(\mathcal{D}\) and any interpretation of the TDG signature
\(\Sigma_{\mathrm{TDG}}\) in \(\mathcal{D}\), there exists a unique functor
\(\tilde{F} : E_{\mathrm{sem}} \to \mathcal{D}\) extending that interpretation.
\end{theorem}

\begin{proof}
This is the standard universal property of the free category on a graph:
objects of \(E_{\mathrm{sem}}\) are objects of \(\Sigma_{\mathrm{TDG}}\), and morphisms
are freely generated paths in the underlying graph. Any graph morphism
from \(\Sigma_{\mathrm{TDG}}\) into \(\mathcal{D}\) extends uniquely to a functor.
\end{proof}

\subsubsection{TDG Grammar}

We present TDG both as a signature and as a context-free grammar
over terms.

\paragraph{Nonterminals.}
We use:
\[
Q \quad\text{for explicit pipelines}, \qquad
P \quad\text{for full pipelines (explicit + tacit suffix)}.
\]

\paragraph{Terminals.}
The terminal symbols are:
\[
\RI,\, \EC,\, \ED,\, \RB,\, \TE,\, \HFD,\,
\BRANCH,\, \PARALLEL,\, \LOOP.
\]

\paragraph{Productions.}
The grammar productions are:
\begin{align*}
Q &::= \RI,\\
Q &::= Q;\EC,\\
Q &::= Q;\ED,\\
Q &::= \BRANCH(Q_1,\dots,Q_k),\\
Q &::= \PARALLEL(Q_1,\dots,Q_k),\\
Q &::= \LOOP(Q_{\mathrm{cond}}, Q_{\mathrm{body}}),\\
P &::= Q;\RB;\TE;\HFD.
\end{align*}
Thus TDG presents \(E_{\mathrm{sem}}\) as an algebra of explicit pipelines
with a fixed tacit suffix.

\subsubsection{Shapes and Shape Category}

\begin{definition}[TDG shape]
TDG shapes are generated by the grammar
\[
S ::= \chain \mid
       \branch(S_1,\dots,S_k) \mid
       \parallel(S_1,\dots,S_k) \mid
       \loopshape(S_{\mathrm{cond}}, S_{\mathrm{body}}).
\]
\end{definition}

\begin{definition}[Shape of a TDG term]
The mapping
\(\shape : \text{TDG-terms} \to \text{Shapes}\) is defined inductively by:
\begin{itemize}
  \item \(\shape(\RI) = \chain;\)
  \item \(\shape(Q;\EC) = \shape(Q)\) and similarly \(\shape(Q;\ED) = \shape(Q)\);
  \item \(\shape(\BRANCH(Q_1,\dots,Q_k))
         = \branch(\shape(Q_1),\dots,\shape(Q_k));\)
  \item analogously for \(\PARALLEL\) and \(\LOOP\).
\end{itemize}
\end{definition}

\begin{definition}[Shape category \(\Scat\)]
The shape category \(\Scat\) has:
\begin{itemize}
  \item objects: shapes \(S\) as above;
  \item morphisms: structural rewiring maps between shapes
        (e.g., regrouping, insertion/removal of trivial chains),
        closed under composition and identities.
\end{itemize}
\end{definition}

\subsubsection{Explicit Reduction System}

\begin{definition}[Explicit reduction]
A binary relation \(\to_E\) on \(\Espace\) is an \emph{explicit reduction system}
if it is a typed rewrite system satisfying:
\begin{itemize}
  \item \textbf{typing:} if \(e : X \to E\) and \(e \to_E e'\), then \(e' : X \to E\);
  \item \textbf{TDG compatibility:} reductions act componentwise on TDG
        combinators (\(\BRANCH\), \(\PARALLEL\), \(\LOOP\));
  \item \textbf{termination:} there is no infinite chain
        \(e_0 \to_E e_1 \to_E e_2 \to_E \dots\);
  \item \textbf{local confluence:} if \(e \to_E e_1\) and \(e \to_E e_2\), then
        \(e_1\) and \(e_2\) have a common descendant.
\end{itemize}
We treat these properties as axioms of the explicit reduction system.
\end{definition}

Under these assumptions, every explicit artifact admits a unique normal
form modulo \(\to_E\).

%======================================================================
\subsection{Navigation and Shapes}
%======================================================================

\subsubsection{Navigation Category and Groupoid}

\begin{definition}[Navigation category \(N_0\)]
The navigation category \(N_0\) has:
\begin{itemize}
  \item objects: shapes \(S\);
  \item generating morphisms: navigation steps
        \(g : S \to S'\), such as \(\TRYMODEL\),
        \(\ADDASSUMPTION\), \(\DROPASSUMPTION\),
        \(\MERGE\), \(\RESTART\), each with a well-defined
        source and target shape;
  \item morphisms: finite composites of such generators.
\end{itemize}
\end{definition}

\begin{definition}[Navigation groupoid \(N\)]
The navigation groupoid \(N\) is obtained by freely adjoining inverses
to morphisms of \(N_0\) and quotienting by the relations
\[
g^{-1} \circ g = \id_S, \qquad
g \circ g^{-1} = \id_{S'},
\]
for each generator \(g : S \to S'\).
Thus every navigation step is reversible up to equality in \(N\).
\end{definition}

\subsubsection{Fibration \(p : E \to S\)}

\begin{definition}[Projection to shapes]
The functor \(p : \Espace \to \Scat\) sends each explicit pipeline
\(e \in \Espace\) to its shape
\(p(e) := \shape(e)\), and each explicit morphism to the corresponding
structural morphism between shapes.
\end{definition}

\begin{definition}[Fibers]
For each shape \(S\), the fiber
\[
E_S := p^{-1}(S)
\]
is the category of explicit morphisms with shape \(S\).
\end{definition}

\begin{definition}[Cartesian liftings and reindexing]
For each navigation morphism \(n : S \to S'\) in \(N\), we assume a
cartesian lifting
\[
n^\ast : E_S \to E_{S'}
\]
which reindexes explicit pipelines along \(n\). This yields an indexed
category
\[
E : N^{\mathrm{op}} \to \mathbf{Cat},
\qquad
S \mapsto E_S,\quad
n \mapsto n^\ast.
\]
\end{definition}

%======================================================================
\subsection{Combined System: Grothendieck Construction}
%======================================================================

\subsubsection{Semidirect Product \texorpdfstring{\(N \ltimes E\)}{N ⋉ E}}

\begin{definition}[Grothendieck construction \(N \ltimes E\)]
Given the indexed category
\(E : N^{\mathrm{op}} \to \mathbf{Cat}\), the Grothendieck
construction \(\int_N E\) (also written \(N \ltimes E\)) has:
\begin{itemize}
  \item objects: pairs \((S, e)\) with \(S \in \Obj(N)\) and \(e \in E_S\);
  \item morphisms \((S, e) \to (S', e')\): pairs \((n, f)\) where
        \(n : S \to S'\) in \(N\) and
        \(f : n^\ast(e) \to e'\) in \(E_{S'}\).
\end{itemize}
Composition is given by
\[
(n_1, f_1) \circ (n_2, f_2)
\;=\;
\bigl(
  n_1 \circ n_2,\,
  f_1 \circ n_1^\ast(f_2)
\bigr).
\]
Informally, we may write elements as \((n, e)\) and composition as
\((n_1, e_1) \circ (n_2, e_2) =
 (n_1 \circ n_2,\, e_1 \circ n_1(e_2))\) whenever the lifted action is defined.
\end{definition}

\subsubsection{Rewrite System on \texorpdfstring{\(N \ltimes E\)}{N ⋉ E}}

\begin{definition}[Combined reduction]
The one-step reduction relation \(\to\) on \(N \ltimes E\) is generated by:
\begin{itemize}
  \item \textbf{explicit step:}
        if \(e \to_E e'\) then \((n, e) \to (n, e')\);
  \item \textbf{navigation step:}
        if \(n \to_N n'\) and \(n'\) is defined on \(e\), then
        \((n, e) \to (n', n'(e))\).
\end{itemize}
We assume:
\begin{enumerate}[label=(A\arabic*)]
  \item \(\to_E\) is terminating and locally confluent;
  \item navigation reductions \(\to_N\) have normal forms in each component
        of \(N\);
  \item the lifted action is functorial and preserves \(\RB\), \(\TE\),
        and \(\HFD\).
\end{enumerate}
These are axioms of the FRFP reduction environment.
\end{definition}

%======================================================================
\subsection{Semantics and Tacit Correctness}
%======================================================================

\subsubsection{Semantic Functor}

\begin{definition}[Semantic bracket]
A \emph{semantic functor} is a functor
\[
\llbracket{-}\rrbracket : N \ltimes E_{\mathrm{sem}} \to \Tspace
\]
assigning to each combined configuration a tacit state.
On explicit pipelines \(e \in E_{\mathrm{sem}}\), we set
\[
\llbracket e \rrbracket := \HFD \circ \TE \circ \RBsharp(e).
\]
\end{definition}

\begin{definition}[Navigation invariance]
The functor \(\llbracket{-}\rrbracket\) is \emph{navigation-invariant up to
preorder} if
\[
\llbracket (n, e) \rrbracket \;\sqsubseteq\; \llbracket e \rrbracket
\]
for all \(n,e\) where the action is defined, and equality holds in the
quotient by the induced equivalence on tacit states in confluent cases.
\end{definition}

\begin{definition}[Observable correctness]
Given a pipeline \(P\), its observable correctness outcome is
\[
\gamma(\llbracket P \rrbracket) \in J.
\]
\end{definition}

%======================================================================
\subsection{Main Theorems}
%======================================================================

\subsubsection{Well-Formedness of TDG Pipelines}

\begin{theorem}[Well-formedness]
\label{thm:well-formed}
Every TDG-generated pipeline \(P\) is a well-typed morphism in \(\Cfull\) of
the form
\[
\RI \;\to\; (\EC/\ED)^{\ast} \;\to\; \RB \;\to\; \TE \;\to\; \HFD.
\]
\end{theorem}

\begin{proof}
By induction on TDG derivations. The base case \(\RI\) initializes an object
in \(E\). The constructors \(\EC\) and \(\ED\) have codomain \(E\) and thus
preserve explicit type. Structural combinators (\(\BRANCH\), \(\PARALLEL\),
\(\LOOP\)) act within the explicit category and preserve codomain \(E\)
componentwise. The production for \(P\) enforces the suffix
\(\RB; \TE; \HFD\), whose types are \(E \to T \to T \to T\).
By ETS there are no morphisms \(T \to E\), so:
\begin{itemize}
  \item no tacit operator can appear before \(\RB\);
  \item no explicit operator can appear after \(\RB\).
\end{itemize}
Thus every derivation yields a well-typed morphism of the stated form.
\end{proof}

\subsubsection{Explicit Completeness of TDG}

\begin{theorem}[Explicit completeness]
\label{thm:explicit-completeness}
TDG generates exactly the admissible explicit morphisms in \(\Espace\).
\end{theorem}

\begin{proof}
By construction, \(E_{\mathrm{sem}}\) is the free category on the TDG
signature, and \(\Ufun\) interprets TDG terms as explicit morphisms in
\(\Espace\). Every explicit morphism in the FRFP setting is built from
\(\RI\), \(\EC\), \(\ED\), and the structural operators, hence arises from
a TDG term; conversely, each TDG term denotes a well-typed explicit
morphism. The image of TDG under \(\Ufun\) is thus exactly the admissible
explicit morphisms.
\end{proof}

\subsubsection{Minimality of Primitive Set}

\begin{theorem}[Minimality]
\label{thm:minimality}
The primitive set
\(\{\RI,\EC,\ED,\RB,\TE,\HFD\}\)
is minimal under the Epistemic Spaces axioms: removing any primitive violates at
least one of ETS, HEG, HEC, RB, CP, or AE.
\end{theorem}

\begin{proof}
Briefly:
\begin{itemize}
  \item \(\RI\) is required to start from \(\varnothing\) and generate explicit
        content; without it there is no canonical initiation of pipelines.
  \item \(\EC\) is required for generative explicit reasoning; without it,
        explicit computation collapses.
  \item \(\ED\) is required for explicit diagnostics without tacit evaluation;
        removing it forces diagnostics either into tacit space (violating AE)
        or to be indistinguishable from \(\EC\) (violating CP and the intended
        modal roles).
  \item \(\RB\) is required for explicit-to-tacit handoff; without it there is
        no admissible morphism \(E \to T\), contradicting the requirement that
        pipelines reach tacit evaluation.
  \item \(\TE\) is required for tacit grounding (HEG); without it correctness
        cannot be evaluated in tacit space.
  \item \(\HFD\) is required for final correctness closure (HEC); without it
        there is no morphism corresponding to final acceptance/rejection.
\end{itemize}
Removing any primitive causes at least one Epistemic Spaces axiom to fail, so the
set is minimal.
\end{proof}

\subsubsection{Epistemic Spaces Initiality}

\begin{definition}[Epistemic Spaces framework]
A \emph{Epistemic Spaces framework} is a tuple
\[
\mathcal{F}_1
 \;=\;
(\mathcal{C}, E_{\mathcal{C}}, T_{\mathcal{C}},
 \RI_{\mathcal{C}}, \EC_{\mathcal{C}}, \ED_{\mathcal{C}},
 \RB_{\mathcal{C}}, \TE_{\mathcal{C}}, \HFD_{\mathcal{C}})
\]
satisfying the Epistemic Spaces axioms.
Morphisms in the category \(\Frm_1\) of Epistemic Spaces frameworks are functors
preserving \(E\), \(T\), and all six primitives.
\end{definition}

\begin{theorem}[Epistemic Spaces initiality]
\label{thm:phase1-initial}
There exists an initial object
\(F^{(1)}_{\mathrm{FRFP}} \in \Frm_1\).
\end{theorem}

\begin{proof}
Construct \(F^{(1)}_{\mathrm{FRFP}}\) as the free category on objects
\(E, T\) and morphisms \(\RI,\EC,\ED,\RB,\TE,\HFD\), modulo the equations
encoding the Epistemic Spaces axioms (ETS, HEG, HEC, RB, CP, AE). Any other
Epistemic Spaces framework interprets these generators in a way that satisfies the
same equations, yielding a unique induced functor
\(F^{(1)}_{\mathrm{FRFP}} \to \mathcal{F}_1\). Thus
\(F^{(1)}_{\mathrm{FRFP}}\) is initial in \(\Frm_1\).
\end{proof}
\paragraph{Remark A.34a (Non-trivial content of Epistemic Spaces initiality).}
The statement that $F^{(1)}_{\mathrm{FRFP}}$ is initial in $\mathrm{Frm}_1$ uses the free-category
construction, but its content is more than formal packaging. The Epistemic Spaces axioms (ETS, HEG,
HEC, RB, CP, AE), together with the minimality result (Theorem~A.32), severely constrain which
primitive roles may appear and how they can be composed. Initiality then says that any epistemic
framework respecting exactly these roles and axioms must interpret RI, EC, ED, RB, TE, HFD in a
way that factors uniquely through the FRFP ontology. In particular, attempts to ``hide'' additional
power in explicit space (for example, by letting EC or ED implicitly perform tacit evaluation) break
the axioms rather than yielding a genuinely different Epistemic Spaces framework.

\subsubsection{Epistemic Algebra Canonicality}

\begin{definition}[Epistemic Algebra architecture]
A \emph{Epistemic Algebra architecture} is a tuple
\[
\mathcal{F}_2
 =
(E,\ S,\ p,\ N,\ (\cdot)^\ast,\ N \ltimes E)
\]
satisfying:
\begin{itemize}
  \item TDG presents \(E\) as a free explicit category;
  \item \(N\) is a navigation groupoid on shapes \(S\);
  \item \(p : E \to S\) is a fibration with cartesian liftings \(n^\ast\);
  \item \(N \ltimes E\) is the Grothendieck construction of the indexed
        category \(E : N^{\mathrm{op}} \to \mathbf{Cat}\);
  \item the boundary structure \(\RB\)--\(\TE\)--\(\HFD\) is preserved
        under navigation.
\end{itemize}
Morphisms in the category of Epistemic Algebra architectures preserve all of this
structure.
\end{definition}

\begin{theorem}[Epistemic Algebra canonicality]
\label{thm:phase2-canonicality}
The FRFP Epistemic Algebra architecture
\(F^{(2)}_{\mathrm{FRFP}}\) is a canonical object in the category of
Epistemic Algebra architectures: for every Epistemic Algebra architecture \(\mathcal{A}\),
there exists a unique structure-preserving morphism
\[
F^{(2)}_{\mathrm{FRFP}} \to \mathcal{A}.
\]
\end{theorem}

\begin{proof}
By the universal properties of:
\begin{itemize}
  \item the free explicit category \(E_{\mathrm{sem}}\);
  \item the groupoid completion constructing \(N\);
  \item the Grothendieck construction forming \(N \ltimes E\).
\end{itemize}
Any Epistemic Algebra architecture satisfying the same axioms must interpret the
FRFP construction in a role-preserving way, inducing a unique morphism
from \(F^{(2)}_{\mathrm{FRFP}}\).
\end{proof}

\subsubsection{Confluence of \texorpdfstring{\(N \ltimes E\)}{N ⋉ E}}

\begin{theorem}[Confluence of \(N \ltimes E\)]
\label{thm:n-e-confluence}
Under assumptions \textup{(A1)}--\textup{(A3)}, the reduction relation
\(\to\) on \(N \ltimes E\) is confluent.
\end{theorem}

\begin{proof}[Proof sketch]
Local peaks \((n, e) \leftarrow (n_0, e_0) \to (n', e')\) decompose into:
\begin{itemize}
  \item explicit--explicit peaks, which join by local confluence of
        \(\to_E\);
  \item navigation--navigation peaks, which join by normalization in each
        connected component of \(N\);
  \item mixed peaks, which join because the lifted action commutes with
        explicit reduction and preserves \(\RB,\TE,\HFD\).
\end{itemize}
Together with termination up to navigation-normalization, this yields
global confluence.
\end{proof}
\paragraph{Remark A.37a (Non-trivial aspects of Epistemic Algebra structure).}
While Theorems~A.36 and A.37 are phrased in terms of standard categorical constructions
(free categories, groupoid completions, Grothendieck constructions), several aspects of the
Epistemic Algebra architecture are mathematically substantive:

\begin{itemize}
  \item \textbf{Compatibility with RB--TE--HFD.} The lifted navigation action and the
  Grothendieck construction are required to preserve the boundary structure. This rules out
  navigation steps that implicitly smuggle tacit evaluation back into explicit space or that
  duplicate/delete the RB--TE--HFD suffix. Showing that the action can be defined so as to
  commute with explicit reduction while preserving RB, TE, HFD is a genuine structural
  constraint on admissible navigation schemes.
  \item \textbf{Confluence of the combined reduction.} Confluence of $\to$ on $N \ltimes E$ is not
  automatic from confluence of $\to_E$ and normalization in $N$. The proof must analyze
  mixed peaks and use the functoriality of the lifted action, together with preservation of
  RB, TE, and HFD, to ensure that explicit and navigation steps can be permuted without
  changing the normal form.
  \item \textbf{Interaction with semantics.} The semantic correctness theorem (Theorem~A.38)
  depends crucially on the fact that navigation-invariant semantics on $N \ltimes E$ respects
  the Interaction Protocols constraints (no illicit $T \to E$ morphisms) and the boundary-preservation
  properties above. This is what forces correctness judgments to depend only on the tacit
  evaluation of the explicit pipeline, not on the particular exploration trajectory.
\end{itemize}

These non-trivial interactions justify treating Epistemic Algebra canonicality as more than a choice of
notation: any alternative architecture satisfying the same constraints must reproduce this
structure and, in particular, cannot avoid the resulting preservation and non-automation results
without breaking one of the axioms.

\subsubsection{Semantic Correctness for Epistemic Structure}

Collecting the constructions above, we obtain a semantic correctness
theorem for Epistemic Structure FRFP pipelines.

\begin{theorem}[Semantic correctness]
\label{thm:semantic-correctness}
Let \(P\) be any TDG-generated pipeline that is Interaction Protocols admissible and
Epistemic Spaces/Algebra well-typed. Then:
\begin{enumerate}
  \item \(P\) has a unique normal form in \(N \ltimes E\) modulo navigation;
  \item the observable correctness outcome \(\gamma(\llbracket P \rrbracket)\)
        is invariant under:
        \begin{itemize}
          \item explicit reduction \(\to_E\),
          \item navigation in \(N\),
          \item composition in the semidirect product \(N \ltimes E\).
        \end{itemize}
\end{enumerate}
In particular, correctness judgments depend only on the (tacitly evaluated)
semantic content of the explicit pipeline, not on the particular exploration
trajectory used to construct it.
\end{theorem}

\begin{proof}[Proof sketch]
Item~(1) follows from Theorem~\ref{thm:n-e-confluence} and the termination
assumptions (A1)--(A3). Item~(2) follows from navigation invariance of the
semantic functor, the Interaction Protocols constraints (which rule out illicit T-to-E
feedback or hidden tacit reconstruction), and the fact that \(\RB,\TE,\HFD\)
are preserved under navigation and do not appear in the explicit phase.
\end{proof}

\subsection{Worked example: sepsis consult as an FRFP pipeline}
\label{sec:cp-sepsis-example}
We sketch a concrete instance of the FRFP calculus corresponding to the sepsis consult scenario
of \S\ref{sec:case-sepsis}. The goal is to show how a real conversational workflow is
represented as a TDG term, reduced to normal form in $N \ltimes E$, and checked for
Interaction Protocols admissibility and Epistemic Spaces/Algebra well-typedness.

\paragraph{TDG and free explicit category.}
Let the TDG have objects
\[
  \Ob(E_0) = \{\textsf{Ctx}, \textsf{Input}, \textsf{CaseSummary},
  \textsf{DraftPlan}, \textsf{CandidatePlan}, \textsf{Report}, \textsf{Decision}\}
\]
and generators
\begin{align*}
  \mathsf{RB} &: \textsf{Ctx} \to \textsf{Input}, \\
  \mathsf{prep} &: \textsf{Input} \to \textsf{CaseSummary}, \\
  \mathsf{ehr} &: \textsf{CaseSummary} \to \textsf{CaseSummary}, \\
  \mathsf{calc} &: \textsf{CaseSummary} \to \textsf{CaseSummary}, \\
  \mathsf{assist}_0 &: \textsf{CaseSummary} \to \textsf{DraftPlan}, \\
  \mathsf{review}_\mathrm{res} &: \textsf{DraftPlan} \to \textsf{CandidatePlan}, \\
  \mathsf{review}_\mathrm{att} &: \textsf{CandidatePlan} \to \textsf{Report}, \\
  \mathsf{HFD} &: \textsf{Report} \to \textsf{Decision}.
\end{align*}
Let $E$ be the free category on this graph, modulo the Epistemic Algebra equations for $\mathsf{TE}$ and
the RB--TE--HFD boundary. We define
\[
  \mathsf{assist} := \mathsf{assist}_0 \circ \mathsf{calc} \circ \mathsf{ehr}
  : \textsf{CaseSummary} \to \textsf{DraftPlan},
\]
and the explicit sepsis pipeline
\[
  e_{\mathrm{sepsis}} :=
  \mathsf{HFD} \circ \mathsf{review}_\mathrm{att} \circ \mathsf{review}_\mathrm{res} \circ
  \mathsf{assist}_0 \circ \mathsf{calc} \circ \mathsf{ehr} \circ \mathsf{prep} \circ \mathsf{RB}
  : \textsf{Ctx} \to \textsf{Decision}.
\]

\paragraph{Reduction and normal form.}
Let $R_E$ be the explicit reduction system on $E$ that contracts $\mathsf{ehr}$ and $\mathsf{calc}$
into $\mathsf{assist}$ inside the $\mathsf{assist}_0$ wrapper:
\[
  \mathsf{assist}_0 \circ \mathsf{calc} \circ \mathsf{ehr} \;\to\; \mathsf{assist}.
\]
By confluence of $R_E$ and compatibility with the RB--TE--HFD boundary, the normal form of
$e_{\mathrm{sepsis}}$ is
\[
  e_{\mathrm{sepsis}}^{\ast} :=
  \mathsf{HFD} \circ \mathsf{review}_\mathrm{att} \circ \mathsf{review}_\mathrm{res} \circ
  \mathsf{assist} \circ \mathsf{prep} \circ \mathsf{RB}.
\]
Lifting to $N \ltimes E$ and applying the confluence theorem for the combined reduction (Epistemic Algebra)
ensures that any admissible navigation of the conversational UI yields a unique normal form
equivalent to $e_{\mathrm{sepsis}}^{\ast}$ up to the RB--TE--HFD suffix.

\paragraph{Interaction Protocols admissibility.}
We interpret the Interaction Protocols axioms IL, AR, MD, and NTER as constraints on realizations of
$e_{\mathrm{sepsis}}$ in a concrete UI:
\begin{itemize}
  \item IL: the only interaction channels between AI and human tacit state are via explicit objects in
  $E$ (case summaries, drafts, candidate plans, reports). No morphisms $T \to E$ are introduced.
  \item AR: there is no morphism in the realization that effects clinical actions without passing
  through $\mathsf{HFD} \circ \mathsf{review}_\mathrm{att}$; the AI cannot place orders or write notes.
  \item MD: all model-dependent computation is represented by $\mathsf{assist}$ and its internal
  decomposition; tacit evaluation (clinician judgment) is localized at $\mathsf{review}_\mathrm{res}$ and
  $\mathsf{review}_\mathrm{att}$.
  \item NTER: the implementation forbids any arrow that would encode a tacit correctness oracle in
  explicit space (no reconstruction of $T$ from model outputs). In particular, the realization of
  $\mathsf{assist}$ is constrained to produce drafts and rationales, not authoritative decisions.
\end{itemize}
Under these conditions, any concrete realization of $e_{\mathrm{sepsis}}$ satisfies the Interaction Protocols
interface axioms and is therefore admissible.

\paragraph{Epistemic Spaces/Algebra well-typedness.}
In Epistemic Spaces, each generator is assigned an epistemic role: $\mathsf{prep}$ and $\mathsf{assist}$ are
explicit-computation arrows; $\mathsf{review}_\mathrm{res}$ and $\mathsf{review}_\mathrm{att}$ are explicit
presentation points at which tacit evaluation in $T$ occurs; $\mathsf{RB}$ and $\mathsf{HFD}$ are
boundary morphisms. The typing rules of the Epistemic Spaces ontology then derive a composite judgment
showing that $e_{\mathrm{sepsis}}^{\ast}$ is a well-typed explicit pipeline with correctness localized in
tacit post-composition at $\mathsf{review}_\mathrm{att}$. In Epistemic Algebra, the same term appears as a
normal form in $N \ltimes E$; compatibility of reduction with the semantics ensures that the
correctness predicate applied to $e_{\mathrm{sepsis}}^{\ast}$ coincides with the protocol-level notion of
FRFP-hallucination (a violation of $t \succeq \mathrm{Req}(e,c)$ at the moment of final acceptance).
Thus the sepsis consult example realizes a concrete conversational UI that is Interaction Protocols admissible,
Epistemic Spaces/Algebra well-typed, and corresponds to a specific normal form in the FRFP calculus.
